\documentclass[10pt,a4paper]{amsart}
\usepackage[margin=19mm]{geometry}
\usepackage{amsmath,amssymb,mathtools}
\usepackage[scr=rsfs,cal=euler]{mathalfa}
\usepackage{xfrac}
\usepackage[unicode,hidelinks,bookmarks=false]{hyperref}
\newcommand{\ie}{i.e.\ }
\newcommand{\cf}{cf.\ }
\newcommand{\resp}{respectively\ }
\newcommand{\Subsection}{\subsection}
\providecommand{\nobreakdash}{\nobreak\hbox{-}}
\setlength{\emergencystretch}{2em}
\multlinegap0pt
\usepackage[shortlabels]{enumitem}
\usepackage{amssymb}
\usepackage{bm}
\usepackage{dsfont}
\usepackage{xcolor}
\usepackage{tikz}
\usepackage{float}
\usepackage{mathtools}
\usepackage{algorithm}\usepackage{algpseudocode}
\newtheorem{thm}{Theorem}
\title{Conflict-free chromatic index of bipartite graphs}
\author{Yuxin Jin, Yuping Gao}
\date{Source: arXiv:2604.24183v1 (CC BY 4.0)}
\begin{document}
\maketitle

\noindent\emph{Source and licence.} Conflict-free chromatic index of bipartite graphs. Version arXiv:2604.24183v1 (CC BY 4.0), distributed by arXiv under the Creative Commons Attribution 4.0 International licence, http://creativecommons.org/licenses/by/4.0/, as recorded in the arXiv licence metadata for this exact version and shown on the abstract page https://arxiv.org/abs/2604.24183v1. Full licence notice: \texttt{licenses/arxiv-2604.24183-CC-BY-4.0.txt}.\par\smallskip

\noindent\textit{Editorial scope.} Counted unit: the article's thm thm:tree (source0.tex lines 168--178). The article's complete original proof of it is delivered with it (source0.tex lines 257--274, one \texttt{proof} environment of 18 lines, which the article places away from the statement). No mathematics is written, repaired or re-derived: the statement and the proof are the article's own lines, in the article's own notation, and no retained line is substituted or re-flowed.  No further source-local context is needed: every cross-reference of every retained line resolves inside the two delivered spans.  Nothing was omitted from any retained span, so the package carries no omission record.  No retained line contains a citation, so the wrapper carries no bibliography and no bibliography entry.  No retained line contains a figure environment, an image-inclusion command or drawing code, so no image is delivered and the package declares no asset dependency.  Editorial formatting only, in addition: an AMS \texttt{amsart} preamble that restates the article's own declarations and macros used by the retained lines, namely the environments \texttt{thm} on separate counters, together with the standard packages those lines need.  The retained environments print in this document as Theorem 1 in order of appearance, whereas the article prints the same items with the numbering of its own full document.  The complete native source of the article as distributed in the arXiv e-print for this exact version is bundled unchanged as \texttt{sources/199217/source0.tex}.
\begin{thm}\label{thm:tree}
Let $T$ be a tree with $|E(T)|\ge 2$. Then $\chi'_{CF}(T)=2$ if and only if
there exists a nonempty proper subset $F\subseteq E(T)$ such that, for every edge
$uv\in E(T)$, one of the following holds:
\begin{enumerate}[label={\rm(\roman*)}]
\item if $uv\in F$, then
$d_F(u)+d_F(v)=2$ or $(d_T(u)-d_F(u))+(d_T(v)-d_F(v))=1$; \label{condition:i}
\item if $uv\notin F$, then
$d_F(u)+d_F(v)=1$ or $(d_T(u)-d_F(u))+(d_T(v)-d_F(v))=2$.\label{condition:ii}
\end{enumerate}
\end{thm}

\begin{proof}[Proof of Theorem~\ref{thm:tree}]
First, assume that $\chi'_{CF}(T)=2$. Let $\varphi$
be a conflict-free edge $2$-coloring of $T$, and let $F:=\{e\in E(T)\colon \varphi(e)=1\}$.
Since both colors are used on $T$, $F$ is a nonempty proper subset of $E(T)$.
Let $uv\in E(T)$ and let $n_i(uv)$ denote the number of edges colored by $i$ in $E_T[uv]$ for $i\in\{1,2\}$.
Since $\varphi$ is conflict-free, either $n_1(uv)=1$ or $n_2(uv)=1$.

If $uv\in F$, then $n_1(uv)=d_F(u)+d_F(v)-1$ and $n_2(uv)=(d_T(u)-d_F(u))+(d_T(v)-d_F(v))$. So either $d_F(u)+d_F(v)=2$ or $(d_T(u)-d_F(u))+(d_T(v)-d_F(v))=1$. Condition~\ref{condition:i} follows.
If $uv\notin F$, then $n_1(uv)=d_F(u)+d_F(v)$ and $n_2(uv)=(d_T(u)-d_F(u))+(d_T(v)-d_F(v))-1$. So either $d_F(u)+d_F(v)=1$ or $(d_T(u)-d_F(u))+(d_T(v)-d_F(v))=2$. Condition~\ref{condition:ii} follows.


Conversely, assume that there exists a nonempty proper subset $F\subseteq E(T)$ satisfying conditions~\ref{condition:i} and~\ref{condition:ii}. Let $\varphi$ be an edge 2-coloring of $T$ such that $\varphi(e)=1$ if $e\in F$ and $\varphi(e)=2$ otherwise.
For any edge $uv\in E(T)$, let $n_i(uv)$ denote the number of edges colored by $i$ in $E_T[uv]$ for $i\in\{1,2\}$.

If $uv\in F$, then either $n_1(uv)=d_F(u)+d_F(v)-1=1$ or $n_2(uv)=(d_T(u)-d_F(u))+(d_T(v)-d_F(v))=1$ by condition~\ref{condition:i}.
If $uv\notin F$, then  either $n_1(uv)=d_F(u)+d_F(v)=1$ or $n_2(uv)=(d_T(u)-d_F(u))+(d_T(v)-d_F(v))-1=1$ by condition~\ref{condition:ii}.
Therefore, all edges of $T$ are satisfied by $\varphi$ and we have $\chi'_{CF}(T)=2$.
\end{proof}

\end{document}
